\documentclass[tikz,border=3mm,multi=tikzpicture]{standalone}
\usepackage{eleve-matematica-ef1}

\begin{document}

% FIGURA 1 — fig-01-leitura-da-regua
\begin{tikzpicture}[matematica figura, x=0.55cm, y=1cm]
  \path[use as bounding box] (-0.40,-0.25) rectangle (14.70,4.20);
  \fill[matemacinzaclaro] (0,0.35) rectangle (14,1.55);
  \draw[matematica contorno] (0,0.35) rectangle (14,1.55);
  \foreach \x in {0,...,14}{
    \draw[matematica eixo] (\x,1.55)--(\x,1.03);
    \node[font=\normalsize,text=matemacinza] at (\x,0.72) {\x};
  }
  \draw[matematica destaque, line width=5pt, rounded corners=2pt] (0,2.65)--(14,2.65);
  \draw[matematica destaque claro, fill=matemalaranjaclaro] (14,2.65)--(13.35,3.00)--(13.35,2.30)--cycle;
  \draw[densely dashed,matematica auxiliar] (0,1.55)--(0,3.35);
  \draw[densely dashed,matematica auxiliar] (14,1.55)--(14,3.35);
  \draw[{Stealth[length=2.5mm]}-{Stealth[length=2.5mm]},matematica contorno] (0,3.55)--(14,3.55);
  \node[matematica valor, fill=white, inner sep=3pt] at (7,3.55) {$14\text{ cm}$};
\end{tikzpicture}

% FIGURA 2 — fig-02-perimetro-de-figuras
\begin{tikzpicture}[matematica figura, x=1cm, y=1cm]
  \path[use as bounding box] (-0.30,-0.25) rectangle (8.70,5.80);
  \fill[matemazulclaro] (1.15,1.15) rectangle (7.15,4.15);
  \draw[matematica destaque, line width=2.6pt] (1.15,1.15) rectangle (7.15,4.15);
  \node[matematica valor, fill=white, inner sep=3pt] at (4.15,4.15) {$6\text{ cm}$};
  \node[matematica valor, fill=white, inner sep=3pt, rotate=90] at (7.15,2.65) {$3\text{ cm}$};
  \node[matematica rotulo] at (4.15,0.45) {somamos os quatro lados};
  \node[matematica resultado] at (4.15,5.15) {contorno: $18\text{ cm}$};
\end{tikzpicture}

% FIGURA 3 — fig-03-leitura-da-jarra
\begin{tikzpicture}[matematica figura, x=1cm, y=1cm]
  \path[use as bounding box] (-0.20,-0.15) rectangle (6.80,7.55);
  \fill[matemazulclaro] (1.20,0.65) rectangle (4.65,4.05);
  \draw[matematica contorno, rounded corners=8pt] (1.20,0.65)--(1.20,6.55)--(4.65,6.55)--(4.65,0.65)--cycle;
  \draw[matematica contorno] (4.65,5.60) .. controls (6.25,5.55) and (6.25,2.25) .. (4.65,2.25);
  \foreach \y/\v in {0.65/0,1.78/200,2.91/400,4.04/600,5.17/800,6.30/1000}{
    \draw[matematica eixo] (0.72,\y)--(1.20,\y);
    \node[font=\normalsize,text=matemacinza,anchor=east] at (0.58,\y) {\v};
  }
  \draw[matematica destaque, line width=2.2pt] (1.20,4.04)--(4.65,4.04);
  \node[matematica resultado] at (3.05,7.05) {$600\text{ mL}$};
\end{tikzpicture}

\end{document}
